Esta interpretación la redactó Claude a partir de los registros de las corridas; no es del autor.

Astra obtuvo \textsc{Certified} en los 15 episodios de las cinco variantes, en el primer intento en 14 de ellos (en $\mu = 0{,}01$, $n = 4$, el primer programa falló con un \texttt{NameError} y el segundo quedó certificado). Hay que leer ese 15 de 15 con dos cautelas.

La primera es la referencia. En estas variantes abiertas $\tau_{\mathrm{ref}}$ no es un óptimo publicado sino la búsqueda propia de 06 por entropía cruzada; donde hay tabla publicada ($f(x_n) - f^*$ con $\mu = 0$), esa búsqueda quedaba a un 1--2\,\% del óptimo para $n \ge 4$. \textsc{Certified} dice, con certificado racional exacto, que el peor caso de los pasos de Astra es estrictamente menor que el de esa búsqueda (mejoras entre 0{,}45\,\% y 1{,}0\,\%); no dice que mejoren la literatura. Los programas de Astra son composiciones recursivas con bloques de paso $\sqrt{2}$ y un paso ``puente'' calculado a partir de las sumas de los bloques, del estilo de las composiciones de \citet{GrimmerShuWang2024}; lo más probable es que sean un redescubrimiento, pero no lo comprobamos contra sus tablas.

La segunda es la memoria. En los episodios 2 y 3 de cada variante Astra recibió, en el resumen de episodios anteriores, el programa ya certificado del episodio 1, y lo volvió a enviar (idéntico o con pasos casi iguales). Son, en rigor, cinco descubrimientos independientes, uno por variante, y diez repeticiones. En la variante con pasos negativos permitidos, Astra usó los mismos cuatro pasos positivos que sin ellos, en otro orden: esta corrida no aporta evidencia de que los pasos negativos ayuden.
